% NEREUS: Modular Context-Aware Probabilistic Ship Trajectory Forecasting
% MIWAI 2026 — The 19th International Conference on Multi-disciplinary Trends
%               in Artificial Intelligence, Halifax, Canada, October 10-12, 2026
%               Pages limit: 12 Pages
%
\documentclass[runningheads]{preamp/llncs}
%
\usepackage[T1]{fontenc}
\input{preamp/mydefs}
\input{preamp/acronyms}

%\renewcommand\UrlFont{\color{blue}\rmfamily}
%
\begin{document}
%
\title{NEREUS: Modular Context-Aware Probabilistic Ship Trajectory Forecasting}
%
\titlerunning{NEREUS: Context-Aware Ship Trajectory Forecasting}
%
% ====================== AUTHOR BLOCK ======================
% ---- CAMERA-READY authors (UNCOMMENT after acceptance) ----
% \author{
%     Ghassan Al-Falouji  \orcidlink{0000-0001-9458-6759} \and
%     Michel Spils        \orcidlink{0000-0002-6431-6085} \and
%     Ben Biesenbach      \orcidlink{0009-0002-3552-4667} \and
%     Sven Tomforde       \orcidlink{0000-0002-5825-8915}
%     }
% %
% %
% \authorrunning{G. Al-Falouji et al.}
% % First names are abbreviated in the running head.
% % If there are more than two authors, 'et al.' is used.
% %
% \institute{
%     Kiel University,~Germany \\ \email{\{gaf,msp,bbi,st\}@informatik.uni-kiel.de}
%     }
%
% ---- BLIND-REVIEW block (DELETE this block after acceptance) ----
\author{Anonymous Author(s)}
\authorrunning{Anonymous}
\institute{Affiliation withheld for double-blind review}
% ==========================================================
%
\maketitle
%
%%%%%%%%%%%%%%%%%%%%%%%%%%%%%%%%%%%%%%%%%%%%%%%%%%%%%%%%%%%%%%
%                                                            %
%                       ABSTRACT                             %
%                                                            %
%%%%%%%%%%%%%%%%%%%%%%%%%%%%%%%%%%%%%%%%%%%%%%%%%%%%%%%%%%%%%%
\begin{abstract}
%__ MOTIVATION ___________________________________________
Accurate short-term trajectory prediction for maritime vessels
is a critical capability for
collision avoidance,
traffic monitoring,
and
autonomous maritime navigation.
%__ CHALLENGES _____________________________________________
Most learning-based approaches
model only one or two contextual factors--%
motion history,
vessel-to-vessel interaction,
or navigational geography--%
leaving
complementary signals unexploited,
and
rarely quantifying how these signals interact.
%__ THE PROPOSED SOLUTION __________________________________
In this paper, we present \acs{NEREUS},
a modular deep learning-based framework
for
probabilistic ship trajectory forecasting
that
systematically integrates five contextual streams:
\begin{enumerate*}[label=(\roman*)]
    \item temporal motion dynamics;
    \item spatial context extracted from geo-referenced nautical charts;
    \item static vessel properties such as ship-type and hull dimensions;
    \item social interactions between nearby vessels; and
    \item long-term navigation behaviour from historical \acs{AIS} data.
    %> The previous point is added below to highlight the bayesian-like prior in forming the post. knowledge
\end{enumerate*}
%
%__ PAPER CONTRIBUTION ___________________________________________
An adaptive fusion mechanism
learns to
weight and combine
all available input streams
before a \acs{MDN} decoder generates $K$-probabilistic trajectory modes.
%
\acs{NEREUS}
further incorporates a data-driven trajectory prior---%
either
ship-type-specific density maps
or
trajectory forecasts from a pretrained transformer
or
\acs{GMM}-clustered route bundles of those forecasts,
to
bias predictions toward common navigation patterns.
%
Because every stream is a switchable module over a shared latent space,
the framework supports a controlled ablation
that measures each signal's contribution under a matched capacity budget.
%
%__ EVALUATION SNEAK-PEEK ___________________________________
Trained on two years of \acs{AIS} data from the Kiel--Flensburg region, Germany,
and
evaluated in-domain,
the resulting ablation shows that enriching the trajectory input
with map and ship-to-ship context cuts the \acs{ADE} from $61.5$ to $46.8$\,m,
and
that the long-term prior is the single strongest contextual module,
while
stacking further modules yields diminishing returns under matched capacity.
%
%
% the framework supports a controlled ablation that
% yields three findings of practical consequence:
% % FINDINGS >>
% % First,
% %(1)
% First,
% enriching the trajectory input with map and ship-to-ship context cuts the \acs{ADE} from $61.5$ to $46.8$\,m.
% %(2)
% Second,
% among contextual modules, the prior-module is the single strongest signal,
% and
% %(3)
% finally,
% stacking social- and prior-modules on top yields diminishing-to-negligible returns
% under matched capacity;
% a simple empirical density prior matches or beats a learned transformer prior.

\keywords{
        ship trajectory prediction
        \and
        AIS
        \and
        graph neural networks
        \and
        mixture density network
        \and
        maritime collision avoidance
        \and
        probabilistic forecasting
        \and
        situational awareness
        }
\end{abstract}
%
%%%%%%%%%%%%%%%%%%%%%%%%%%%%%%%%%%%%%%%%%%%%%%%%%%%%%%%%%%%%%%
%                                                            %
%                       INTRODUCTION                         %
%                                                            %
%%%%%%%%%%%%%%%%%%%%%%%%%%%%%%%%%%%%%%%%%%%%%%%%%%%%%%%%%%%%%%
\section{Introduction}
\label{sec:intro}

%__ MOTIVATION ___________________________________
Maritime transport accounts for over 80\,\% of global trade by volume,
making navigation
safety and efficiency
on the world's waterways a matter of
economic and humanitarian importance~\cite{development2023review}.
%
The \ac{AIS},
mandated for all commercial vessels above a certain size,
continuously broadcasts each ship's
position, speed, heading, and identity.
% at one-second to several-minute intervals.
%
% This stream of trajectory data presents a unique opportunity:
If an
autonomous system can accurately forecast where vessels will be at least over the next minutes, it can
detect imminent conflicts
and
support collision-avoidance decision-making
long before human operators would react.

%__ CHALLENGES _________________________________________
Despite apparent similarity to pedestrian and vehicle trajectory forecasting,
the maritime domain poses several distinctive challenges.
\begin{enumerate*}[label=(\roman*)]
    \item Vessels have long stopping distances and turning radii,
    so small angular errors over a five-minute horizon translate into hundreds of metres.

    \item The navigational environment is richly structured:
    coastlines, harbour entrances, shallow-water zones, and traffic separation schemes
    impose hard constraints that data-driven open-water models ignore.

    \item The interaction graph is sparse yet safety-critical---%
    one approaching vessel can matter more than several distant ones---%
    and
    risk is conventionally quantified through \ac{CPA} indicators~\cite{imo1972convention}.

    \item Maritime trajectories are often genuinely multi-modal.
    For instance,
    a vessel approaching a harbour entrance may turn into any of several berthing lanes,
    so a predictor must represent uncertainty rather than average over it.
\end{enumerate*}

%__ LIMITATIONS OF CURRENT SOLUTIONS _________________________________
Existing deep predictors address these challenges only in isolation, and rarely isolate how much each context source contributes (\Cref{sec:background})
% Existing deep predictors address these challenges only in isolation:
% %
% recurrent encoders capture motion but ignore geography and neighbours;
% %
% transformers improve long-range paths but model neither interaction nor real-valued uncertainty;
% %
% and
% recent maritime spatio-temporal graphs add interaction
% yet
% seldom encode the navigational geometry of the chart.
% %
% Crucially,
% they rarely isolate how much each context source actually contributes.
%
% Existing deep predictors address subsets of these challenges in isolation.
% %
% Recurrent \ac{GRU}/\ac{LSTM} encoders
% capture motion but ignore geography and neighbours~\cite{murray2021ais,capobianco2021deep}.
% %
% TrAISformer~\cite{nguyen2024transformer}
% improves long-range paths over discretised cells
% but models neither interaction nor real-valued uncertainty.
% %
% Recent maritime spatio-temporal graphs model interaction
% and even long-term patterns~\cite{feng2022stgcnn,zhang2024learning,huang2024vessel},
% yet they seldom incorporate the navigational geometry of the chart,
% and
% rarely isolate how much each context source actually contributes.
%

%__ SCOPE & SOLUTION PROVIDED IN THIS PAPER _______________________________
This paper presents \ac{NEREUS}, a
modular,
context-aware,
probabilistic framework
for short-term ship trajectory forecasting.
%
It is organised as a Bayesian-inspired decomposition---%
prior (long-term route structure),
likelihood (map and social context),
and
posterior (short-term kinematics)---%
so that each contextual signal is a switchable module.
%
We use this design to measure each contextual signal's contribution
under a matched-capacity budget.


% CODE >>
Code for the \ac{AIS} processing pipeline\footnote{\url{https://anonymous.4open.science/r/ais-processing-26miwai/}}
and
the \ac{NEREUS} framework\footnote{\url{https://anonymous.4open.science/r/nereus-26miwai/}}
is made publicly available.


The key contributions are:
%
\begin{enumerate}
    \item \textbf{A unified probabilistic formulation
    for context-aware maritime forecasting.}
    Integrating
    a long-term prior over route structure,
    a contextual likelihood over social interaction evidence
    within the environmental constraints,
    and
    a short-term probabilistic prediction posterior.

    \item \textbf{Maritime-specific graph representation.}
    We construct interaction graphs whose edge attributes encode
    the regulatory risk geometry underlying the \ac{COLREGs},
    exposing encounter-risk and right-of-way geometry
    to the interaction encoder
    rather than relying on raw proximity.

    \item \textbf{Data-driven trajectory priors.}
    Integrating the knowledge of ``usual-navigation`` long-term patterns,
    through three comparable realisations:
    (a)~ship-type-specific density maps capturing historical route distributions per vessel category,
    (b)~ a frozen pretrained TrAISformer~\cite{nguyen2024transformer} repurposed as a spatial-occupancy prior,
    and
    (c)~a \ac{GMM} over its latent space yielding behaviour archetypes.
    This prior is consumed by an attention-conditioned \ac{CNN}
    that queries the nautical scene map,
    allowing historical patterns to guide where on the chart the model should focus.
\end{enumerate}

%The remainder of this paper is organised as follows.
%\Cref{sec:background} reviews related work.
%\Cref{sec:model} describes the \ac{NEREUS} architecture in detail.
%\Cref{sec:experiments} and~\Cref{sec:results} present the experimental
%setup and results.
%\Cref{sec:conclusion} concludes.

%%%%%%%%%%%%%%%%%%%%%%%%%%%%%%%%%%%%%%%%%%%%%%%%%%%%%%%%%%%%%%
%                                                            %
%                        BACKGROUND                          %
%                                                            %
%%%%%%%%%%%%%%%%%%%%%%%%%%%%%%%%%%%%%%%%%%%%%%%%%%%%%%%%%%%%%%
\section{Background}
\label{sec:background}

%___ INTRO: Taxonomy _______________________________
Vessel trajectory prediction
has evolved along four largely complementary axes:
% (1) >
sequence modelling of \ac{AIS} streams,
% (2) >
interaction modelling between vessels,
% (3) >
environmental and chart-based context,
and
% (4) >
probabilistic forecasting with long-term route priors.
%
Existing maritime methods typically advance one or two of these axes at a time,
leaving the others only weakly represented.
\ac{NEREUS} addresses the collected gaps (see~\Cref{tab:taxonomy}).


%__ (1) Sequence Modelling __________________________
% \paragraph{Sequence models for \ac{AIS} trajectories.}
Early data-driven predictors modelled \ac{AIS} trajectories as temporal sequences of position,
speed, and course, using recurrent encoders
such as \ac{LSTM}, \ac{GRU}, and sequence-to-sequence variants~\cite{murray2021ais,capobianco2021deep,er2025navigating}.
%
More recent work has adopted transformers,
either over discretised spatial cells~\cite{nguyen2024transformer}
or
in hybrid recurrent-attention architectures
to better capture long-range dependencies~\cite{jiang2023trfm,xue2024g}.
%
These models are effective at short-term motion prediction~\cite{li2023aisreview},
yet comparative evaluations also show
that increasingly complex sequence encoders do not necessarily translate into clear gains
over simpler fusion designs~\cite{slaughter2025vessel}.
% Limitation >
Their common limitation is structural:
vessels are usually modelled in isolation,
without explicit interaction reasoning or navigational context,
and
the output is often a single deterministic future.

%___ (2) Interaction Aware Modelling ____________________
% \paragraph{Interaction-aware and graph models.}
To move beyond single-vessel forecasting,
trajectory prediction increasingly models vessel motion
as socially and spatially coupled.
%
The underlying idea follows earlier work on pedestrians and road agents,
where social pooling~\cite{alahi2016social,lee2017desire}
and interaction graphs~\cite{li2019grip} were introduced to capture mutual influence.
%
Maritime extensions use spatio-temporal graph convolution and attention mechanisms
to represent nearby traffic~\cite{feng2022stgcnn,wang2024vessel,huang2024vessel};
more recent variants
jointly model dynamic interactions and longer temporal structure~\cite{zhang2024learning},
while
generative approaches have also been used to learn collision-avoidance behaviour
beyond explicitly coded rules~\cite{jia2023ragan}.
% Limitation >
Despite these advances,
interaction is still commonly approximated through Euclidean neighbourhood graphs
rather than \ac{COLREGs}-relevant risk geometry,
is often encoded only at a single time step or over a largely static graph,
and
is rarely integrated with the geometry of the nautical chart itself.

%___ (3) Charts and Environment aware Modelling ___________
% \paragraph{Environmental and map context.}
A separate line of work
incorporates the environment in which motion unfolds.
%
In autonomous driving,
rasterised or vectorised maps provide scene structure
for trajectory prediction~\cite{lee2017desire,gao2020vectornet},
and
heatmap-based outputs can align predictions with spatial priors~\cite{gilles2022gohome}.
%
Maritime forecasting
has begun to follow this direction by incorporating
weather, currents, and other oceanographic variables,
often fused with vessel interactions
in a graph-based model~\cite{liu2026multimodal,slaughter2025vessel}.
% Limitation >
Yet the most immediate navigational constraints of the maritime domain remain underused:
chart geometry such as coastlines, shallow water, restricted areas, or ferry lanes
is seldom encoded explicitly.
%
When map context is included,
it is also often tied to absolute-position crops,
which risks learning region-specific layouts rather than transferable navigational constraints.


%___ (4) Long Term Vessel Routes _______________________
% \paragraph{Probabilistic output and long-term priors.}
A further distinction concerns how future uncertainty is represented.
Multi-modal prediction
is commonly addressed through
mixture density networks~\cite{bishop1994mixture},
conditional variational autoencoders~\cite{lee2017desire,ivanovic2019trajectron,salzmann2020trajectron++},
and
generative adversarial models~\cite{jia2023ragan};
%
recent maritime work has also started to quantify uncertainty
for collision estimation more explicitly~\cite{jurkus2025enhancing}.
%
Long-term route structure, by contrast,
is typically injected through
clustering,
historical density maps,
or
heatmap priors~\cite{nguyen2024transformer}.
% Limitation >
In maritime prediction, however,
probabilistic forecasting is still the exception rather than the norm,
and models that do represent uncertainty rarely combine an explicit long-term route prior
with a short-term interaction-aware posterior.
%
Safety evaluation is likewise often reduced
to distance-based thresholds, rather than metrics that reflect
navigational rules and encounter geometry.


%___ Framing NEREUS vs SOTA ____________________
% \paragraph{Positioning of \ac{NEREUS}.}

\ac{NEREUS}
treats these four axes as complementary rather than competing
and
unifies them in a single probabilistic maritime forecaster:
interaction-aware sequence modelling,
explicit rasterised chart geometry,
and
long-term prior,
with
vessel interactions defined by \ac{COLREGs}-relevant encounter geometry
rather than proximity alone.
%
Unlike monolithic architectures that entangle all context sources,
it keeps each component modular,
so their individual and joint contributions can be measured
under a matched-capacity budget,
thus attributing their individual performance.

% \ac{NEREUS}
% is designed around the observation that these four axes are complementary
% rather than competing.
% %
% It combines
% interaction-aware sequence modelling,
% explicit rasterised chart geometry,
% and
% a switchable long-term prior within a probabilistic maritime forecaster.
% %
% Instead of defining vessel interactions through proximity alone,
% it constructs edges from \ac{COLREGs}-relevant encounter features
% such as \ac{TCPA}, \ac{DCPA}, relative bearing, and risk indicators.
% %
% Unlike monolithic architectures that entangle all context sources,
% \ac{NEREUS} keeps each component modular,
% allowing their individual and joint contributions to be measured
% under a matched-capacity budget.
% %
% This modular design makes it possible to attribute performance
% to individual context srouces,
% rather than to improve aggregate forecasting accuracy.


\begin{table}[t]
% \caption{Where representative methods sit on the four axes identified above
% (\checkmark~present, $\circ$~partial/coarse, --~absent). Every row is discussed
% in this section or used as a baseline; \ac{NEREUS} is the only maritime
% forecaster that combines all four axes and evaluates each under a controlled,
% matched-capacity ablation.}
\caption{Coverage of the four axes by representative methods (\checkmark present, $\circ$ partial, -- absent).}

\label{tab:taxonomy}
\centering\small
\setlength{\tabcolsep}{5pt}
\begin{tabular}{llcccc}
\toprule
Method & Domain & Multi-modal & Interaction & Chart map & Long-term prior \\
\midrule
GRU/LSTM seq2seq~\cite{murray2021ais,capobianco2021deep} & maritime & -- & -- & -- & -- \\
TrAISformer~\cite{nguyen2024transformer}          & maritime & $\circ$ & -- & -- & $\circ$ \\
IS-STGCNN~\cite{feng2022stgcnn}                 & maritime & -- & \checkmark & -- & -- \\
DSTGAN~\cite{zhang2024learning}                      & maritime & $\circ$ & \checkmark & -- & \checkmark \\
DualSTMA~\cite{huang2024vessel}                 & maritime & -- & \checkmark & -- & -- \\
RAGAN~\cite{jia2023ragan}                          & maritime & \checkmark & \checkmark & -- & -- \\
DESIRE~\cite{lee2017desire}                               & driving  & \checkmark & \checkmark & $\circ$ & -- \\
VectorNet~\cite{gao2020vectornet}/GOHOME~\cite{gilles2022gohome}    & driving  & $\circ$ & \checkmark & \checkmark & $\circ$ \\
\midrule
\textbf{NEREUS (ours)}                             & maritime & \checkmark & \checkmark & \checkmark & \checkmark \\
\bottomrule
\end{tabular}
\end{table}

%%%%%%%%%%%%%%%%%%%%%%%%%%%%%%%%%%%%%%%%%%%%%%%%%%%%%%%%%%%%%%%
%                                                             %
%                        THE FRAMEWORK                        %
%                                                             %
%%%%%%%%%%%%%%%%%%%%%%%%%%%%%%%%%%%%%%%%%%%%%%%%%%%%%%%%%%%%%%%
\section{The NEREUS Framework}
\label{sec:model}

\paragraph{Problem setup.}
At reference time $t_0$,
let $\mathcal{V}_{t_0}=\{v_1,\dots,v_N\}$
denote the vessels present in the scene.
For each vessel $v_i$,
the observed trajectory is $\mathbf{X}_i=(\mathbf{x}_i^{t_0-L+1},\dots,\mathbf{x}_i^{t_0})$
of length $L$, see \Cref{app:features} for details.
%where $\mathbf{x}_i^t$ contains
%relative east/north displacement, speed, course, acceleration, and turn rate (see \Cref{app:features}).
%
The absolute \acs{UTM} position $\mathbf{p}_i^{t_0}$ is retained separately
for chart sampling and evaluation.
%
Static vessel properties,
including ship type and hull geometry, are represented by $\mathbf{s}_i\in\mathbb{R}^{8}$.
%
For vessel pairs within distance $d_{\max}$,
we construct directed edges with attributes $\mathbf{e}_{ij}\in\mathbb{R}^{25}$
%
encoding \ac{COLREGs}-relevant encounter geometry
% (\ac{TCPA}, \ac{DCPA}, relative distance and speed, course difference, and relative bearing)
together with the static features of both endpoints;
the full field list is given in \Cref{tab:features}.

\begin{table}[H]
    \caption{NEREUS input features by group, in code order.
    Node dim 8; static dim 8; edge dim 25; four map channels; single-channel prior.}
    \label{tab:features}
    \centering\small
    \setlength{\tabcolsep}{6pt}
    \begin{tabular}{p{2.5cm}p{8.5cm}}
        \toprule
        Group & Features (per time step unless static), in order \\
        \midrule
        Trajectory (8) & $\Delta x$; $\Delta y$; SOG; acceleration;
        $\sin$/$\cos$ course; $\sin$/$\cos$ turn rate \\
        Static (8) & hull offsets \{to\_bow, to\_stern, to\_port, to\_starboard\};
        ship-type one-hot \{sailing, cargo, passenger, other\} \\
        Interaction (25) & distance; relative SOG; \acs{TCPA}; \acs{DCPA};
        collision-risk; $\sin$/$\cos$ course-difference; $\sin$/$\cos$ relative bearing;
        ego static (8); neighbour static (8) \\
        Map (4) & distance to coastline; distance to restricted area; distance to ferry
        route; water depth ($500\,$m crop, 50\,m grid) \\
        Prior (1) & ship-type density map \emph{or} TrAISformer-path heat-map \emph{or}
        \acs{GMM} cluster-density mixture, calculated at runtime \\
        \bottomrule
    \end{tabular}
\end{table}

% comprising
% relative distance and speed,
% \ac{TCPA}, \ac{DCPA}, collision risk,
% sine/cosine encodings of course difference and relative bearing,
% and
% the static features of both endpoints.
% \gaf{THere is a pretty large overlap between the NEREUS Framework text and the table with data at the end, maybe we can shorten there.}
% %
% Given an ego vessel $v_1$,
% the goal is to predict its next $T$ steps
% as a $K$-component Gaussian mixture over relative displacements,
% which is integrated to absolute positions for evaluation.


% \paragraph{Overview.}
\ac{NEREUS} follows a modular encode--fuse--decode design (\Cref{fig:architecture}).
%
A shared trajectory encoder first produces vessel-level motion embeddings.
Optional context modules then enrich the ego representation
with social interaction, chart context, and a long-term route prior.
Each active module is projected to a common latent space of dimension $H{=}256$,
fused through a learnable softmax-weighted sum,
and passed to an autoregressive \ac{MDN} decoder that predicts a multi-modal future trajectory.

\begin{figure}[tb]
  \centering
  \includegraphics[width=.9\linewidth]{images/nereus_dataflow-1.pdf}
  \caption{\ac{NEREUS} architecture. Five contextual data streams are fused via
    learnable weights before the \ac{MDN} decoder. Map Module, Social Module and Prior are optional.}
  \label{fig:architecture}
\end{figure}


\paragraph{Trajectory encoder and static embedding.}
A shared single-layer \ac{GRU} encodes each vessel history,
$\mathbf{h}_i^{\text{enc}}=\text{GRU}(\mathbf{X}_i)$,
followed by a projection to the common latent dimension $H$;
this trajectory encoder is always active.
%
In parallel,
a linear layer embeds the static descriptor $\mathbf{s}_i$,
consisting of one-hot ship type and normalised hull geometry,
so that downstream modules can condition on vessel-specific manoeuvrability.


\paragraph{Social interaction module.}
\ac{NEREUS} supports two interchangeable interaction mechanisms.
The first is a four-head \ac{GAT}~\cite{velivckovic2018graph} applied to layer-normalised vessel encodings,
where
attention over neighbouring vessels is conditioned on the edge attributes $\mathbf{e}_{ij}$.
This makes interaction reasoning explicitly sensitive to \ac{COLREGs}-relevant encounter geometry.
%
The second is \textbf{Social Pooling}~\cite{lee2017desire},
which aggregates neighbouring hidden states in a polar occupancy grid centred on the ego vessel.
Unlike the \ac{GAT} variant,
this pooling operator does not consume explicit edge attributes;
distance and bearing are only represented through grid placement,
while quantities such as relative speed or \ac{TCPA}/\ac{DCPA} must be inferred
implicitly by later layers.

\paragraph{Map context module.}
To encode navigational constraints,
we rasterise a geo-referenced chart into a four-channel scene tensor containing
distance to coastline, distance to restricted area, distance to ferry route, and water depth.
%
A $500\,\text{m}\times500\,\text{m}$ patch centred at the ego position
$\mathbf{p}_1^{t_0}$ is extracted and encoded by a \textbf{ScenePoolingCNN}.
%
As an alternative,
\textbf{MapAttention} treats the prior representation as a query over CNN feature tokens,
allowing the route prior to focus the model on chart regions that are most relevant
to the predicted motion.

\paragraph{Prior module.}
\ac{NEREUS} can additionally incorporate one of three optional long-term route priors.
% 1:
\textbf{Density}
uses ship-type-specific historical occupancy maps selected according to the ego-vessel type.
% 2:
\textbf{TrAISformer-path}
uses a frozen, pretrained TrAISformer~\cite{nguyen2024transformer},
whose predicted future occupancy is rasterised into a prior map.
% 3:
\textbf{GMM-cluster} models the TrAISformer latent representation $\mathbf{z}$ with a \ac{GMM} ($K{=}16$),
producing route-archetype maps $\{\mathbf{M}_k\}$ that are combined into a prior map
$\sum_k p(k\mid\mathbf{z})\,\mathbf{M}_k$.


\paragraph{Fusion and \ac{MDN} decoder.}
Let $\mathbf{h}^{(m)}$ denote the output of the $m$-th active module.
Each module output is projected to $\mathbb{R}^H$, passed through $\tanh$, and fused by
a learnable, input-independent gate:
%
\begin{equation}
\mathbf{h}=\sum_{m=1}^{M} w_m \tanh\,\big(\mathbf{h}^{(m)}\big),
\qquad
\mathbf{w}=\mathrm{softmax}(\tilde{\mathbf{w}}), \label{eq:fusion}
\end{equation}
%
where $\tilde{\mathbf{w}}$ are free parameters learned jointly with the model.
%
The fused state initialises a single-layer \ac{GRU} decoder that, at each future step $t$, emits the parameters of
a $K$-component Gaussian mixture over the relative displacement $\Delta\mathbf{y}^{t}\in\mathbb{R}^{2}$:
\begin{equation}
\label{eq:mdn}
	p\!\left(\Delta\mathbf{y}^{t}\right)=\sum_{k=1}^{K}\pi_k^{t}\,
	\mathcal{N}\,\big(\Delta\mathbf{y}^{t};\,\boldsymbol{\mu}_k^{t},\,\mathrm{diag}(\boldsymbol{\sigma}_k^{t})^{2}\big),
	\qquad \textstyle\sum_k \pi_k^{t}=1,
\end{equation}
%
with mixture weights $\pi_k^{t}$, means $\boldsymbol{\mu}_k^{t}$, and per-axis standard deviations
$\boldsymbol{\sigma}_k^{t}$. Training minimises the mean per-step negative log-likelihood
$\mathcal{L}=-\frac{1}{T}\sum_{t=1}^{T}\log p(\Delta\mathbf{y}^{t})$ of the ground-truth displacement.
Unless stated otherwise we use $K{=}3$ mixture components and report both a single-hypothesis error
(\acs{ADE}/\acs{FDE} on the mixture mean) and a best-of-$K$ error (\acs{mADE}/\acs{mFDE}); see
\Cref{sec:experiments}.

%%%%%%%%%%%%%%%%%%%%%%%%%%%%%%%%%%%%%%%%%%%%%%%%%%%%%%%%%%%%%%%
%                                                             %
%                         EXPERIMENTS                         %
%                                                             %
%%%%%%%%%%%%%%%%%%%%%%%%%%%%%%%%%%%%%%%%%%%%%%%%%%%%%%%%%%%%%%%

\section{Experimental Setup}
\label{sec:experiments}

\subsection{Dataset}

We use \ac{AIS} records collected from a region in the Baltic sea.
The \textbf{FH10} dataset covers the Kiel Fjord and adjacent
harbour areas.
The dataset is split by day into 70\% training, 10\% validation and 20\% test data.
%The \textbf{DMA10} dataset from the \ac{DMA} covers
%nearby but geographically distinct waters and serves as a
%cross-region test set.
Trajectories are resampled to a fixed 10-second interval
(6 steps per minute) and segmented into windows of 60 observed steps
(10\,min observation) with a 30-step prediction horizon (5\,min).
Nearby vessels within a 500\,m radius form the interaction graph.
In total, the dataset comprises $20\,094\,029$ observation--prediction windows (samples).

\subsection{Baselines}

We compare \ac{NEREUS} against two baselines:
(1)~\textbf{GRU}: a plain encoder--decoder recurrent model with
deterministic output and
(2)~\textbf{DESIRE}~\cite{lee2017desire}: the social pooling with \ac{IOC} refinement predictor adapted to the \ac{AIS} domain.
%% NOTE >>
%	IOC (Inverse Optimal Control) is the control-theoretic sibling of IRL:
%	instead of assuming a known cost/reward and optimising behaviour,
%	you observe expert demonstrations (here, real vessel/agent trajectories)
%	and recover the cost function that makes those demonstrations look near-optimal.
%	You then use that learned cost to score or refine candidate trajectories.
%	In the DESIRE context specifically,
%	IOC is the second stage of the architecture.
%	The CVAE first samples a diverse set of future trajectory hypotheses;
%	the IOC-based ranking-and-refinement module (a recurrent net)
%	then assigns an accumulated long-term reward to each hypothesis
%	and iteratively regresses displacement corrections, pushing the samples toward
%	scene- and interaction-consistent futures.
%	In NEREUS, the 2 IOC refinement iterations hyperparameter is exactly that inner refinement loop unrolled twice.
%(3)~\textbf{TrAISformer}: the transformer-based path forecaster
%re-implemented and trained on the same data split.

\subsection{Training setup \& Hyperparameters}

Training uses AdamW with weight decay $10^{-5}$, a batch size of 1024, and a maximum of 30 epochs with early stopping on the validation loss.
Additionally, a budget of at most 10 hours was set for each model; this budget was reached by each NEREUS module. The learning rate is warmed up linearly over the first epoch and afterwards halved by a plateau scheduler (patience 5, minimum $10^{-6}$).

The GRU baseline is an encoder–decoder with hidden size 256 that autoregressively estimates future positions with a deterministic linear output head, trained with a learning rate of $10^{-3}$.
DESIRE uses a hidden size of 256, a CVAE latent dimension of 16, $K{=}3$ sampled hypotheses, and 2 IOC refinement iterations; social context is aggregated over a log-polar grid (6 rings × 6 wedges, $r_{\min}{=}5$ m) within a 500 m interaction radius, and the scene context CNN maps 4 input channels to 16 feature channels. It is trained with a learning rate of $3\cdot10^{-4}$.

NEREUS shares the GRU backbone and adds the ablated modules:
the social module (GAT with hidden size 64 and 4 attention heads, or social pooling),
the map module (CNN or attention) encodes a 4-channel local map into 128 features.
%
The prior module (density, path, or cluster) likewise outputs a 128-dimensional encoding.
%
In total this yields 30~valid module combinations:
3~social $\times$ 4~prior settings without a map,
plus
3~social $\times$ 3~prior $\times$ 2~map variants,
since the map module requires an active prior.
%
% Table~\ref{tab:nereus_combinations} lists the 30 valid module combinations
% (3 social × 4 prior settings without a map,
% plus 3 social × 3 prior × 2 map variants,
% since the map module requires an active prior).

%\subsection{Ablation Configurations}

% We compare all possible submodule combinations with each other.
% by switching modules on/off. Not all module combinations are possible, since the Map Module depends on the Prior Module.
% The naming convention directly encodes active modules:
% \texttt{base} (encoder + static only),
% \texttt{velivckovic2018graph}/\texttt{pool} (+ social),
% \texttt{cnn} (+ map CNN),
% \texttt{density}/\texttt{cluster}/\texttt{path} (+ prior type),
% and combinations thereof.


% \begin{table}[tb]
% \centering
% \caption{All tested NEREUS module combinations. The encoding GRU and Static projection are always active.}
% \label{tab:nereus_combinations}
% \begin{tabular}{l ccc ccc ccc}
% \toprule
% \textbf{Social} & \multicolumn{3}{c}{None} & \multicolumn{3}{c}{GAT} & \multicolumn{3}{c}{Pool} \\
% \cmidrule(lr){2-4} \cmidrule(lr){5-7} \cmidrule(lr){8-10}
% \textbf{Prior} & None & CNN & Atte & None & CNN & Atte & None & CNN & Atte \\
% \midrule
% --      & \checkmark & --         & --         & \checkmark & --         & --         & \checkmark & --         & --         \\
% Density & \checkmark & \checkmark & \checkmark & \checkmark & \checkmark & \checkmark & \checkmark & \checkmark & \checkmark \\
% Path    & \checkmark & \checkmark & \checkmark & \checkmark & \checkmark & \checkmark & \checkmark & \checkmark & \checkmark \\
% Cluster & \checkmark & \checkmark & \checkmark & \checkmark & \checkmark & \checkmark & \checkmark & \checkmark & \checkmark \\
% \bottomrule
% \end{tabular}
% \end{table}

\subsection{Evaluation Metrics}

% We report:
% \textbf{ADE} (average displacement error over all predicted steps),
% \textbf{FDE@$t$} (displacement error at $t \in \{1,3,5\}$\,min),
% \gaf{discussion: do we actually want to report minK ? seems a bit like cheating.}
% \textbf{minADE} and \textbf{minFDE@$t$} (best-of-$K$ variants for
% multi-modal evaluation).
We report \ac{ADE}- and \ac{FDE}-at $t \in \{1,3,5\}$\,min
on the single-hypothesis mixture mean.
%
The average travelled distance at these steps in 100\,m, 328\,m and 567\,m respectively,
with sailing boats significantly lowering this average.
%
Because the decoder is multi-modal,
we additionally report the best-of-$K$ variants \ac{mADE} and \ac{mFDE}@$t$, which take the mode closest to the ground truth. Best-of-$K$
is intentionally optimistic--it assumes an oracle mode selector--so we treat the
mean-prediction \ac{ADE}/\ac{FDE} as the primary metric and use \ac{mADE}/\ac{mFDE}
only as an upper bound on the multi-modal quality the mixture can attain.

% and three \textbf{collision-risk} metrics computed from predicted hull
% trajectories using a separating-axis-theorem hull distance function:
% mean predicted collision risk, mean minimum predicted inter-vessel
% distance, and collision ratio (fraction of trajectory pairs whose
% predicted hulls overlap).


%%%%%%%%%%%%%%%%%%%%%%%%%%%%%%%%%%%%%%%%%%%%%%%%%%%%%%%%%%%%%%%
%                                                             %
%                           RESULTS                           %
%                                                             %
%%%%%%%%%%%%%%%%%%%%%%%%%%%%%%%%%%%%%%%%%%%%%%%%%%%%%%%%%%%%%%%
\section{Results}
\label{sec:results}


% Table \ref{tab:nereus-metrics} shows the key metrics for the baseline models, the worst and the best evaluated NEREUS models combinations. Even with only the addition of static ship data and a MDN-Head instead of a linear one decreases the ADE from 65.53m to 61.45m.
% Figure~\ref{fig:example trajectories} contains several trajectories from the test set. Frequently used path, for example into a berthing area are well predicted, even when they contain curves. Other trajectories contain completely unexpected turns.
\Cref{tab:nereus-metrics}
reports the key metrics for the baselines
and for the worst and best \ac{NEREUS} configurations.
Even adding only static ship attributes
and an \ac{MDN} head in place of a deterministic linear head
lowers the \ac{ADE} from $65.53$ to $61.45$\,m.
%\gaf{added a sentence for the minFDE values}
Compared to the DESIRE approach the multimodality is more pronounced. $\text{NEREUS}_{\text{best}}$ has a similar or even worse FDE, but considers a wider variety of possible and likely trajectories, leading to a lowered minFDE.

\Cref{fig:example_trajectories} shows several test-set trajectories.
approaches into a berthing area,
are predicted well even when they contain curves,
whereas
some trajectories contain genuinely unexpected turns that the model cannot anticipate.



% >>> Sec.~4.2 (its social-pooling mechanism is already covered by the NEREUS 'Pool' variant).

\begin{table}[htb]
\centering
\caption{Key metrics for baseline models and the worst and best NEREUS configurations. $\text{NEREUS}_{\text{worst}}$ has all modules turned off; $\text{NEREUS}_{\text{best}}$ has the optional modules GAT, attention map, and density-based prior.}
\label{tab:nereus-metrics}
\begin{tabular}{lccccccc}
\toprule
& & \multicolumn{2}{c}{$t=1$} & \multicolumn{2}{c}{$t=3$} & \multicolumn{2}{c}{$t=5$} \\
\cmidrule(lr){3-4} \cmidrule(lr){5-6} \cmidrule(lr){7-8}
Model & ADE[m] & FDE[m] & minFDE[m] & FDE[m] & minFDE[m] & FDE[m] & minFDE[m] \\
\midrule
GRU                             & 65.53 & 16.73 & -     & 78.87  & -      & 160.87 & -      \\
DESIRE                          & 61.12 & 14.81 & 14.76 & 74.02  & 73.83  & 150.30 & 149.96 \\
$\text{NEREUS}_{\text{worst}}$  & 61.45 & 13.67 & 10.76 & 74.21  & 54.08  & 154.56 & 109.16 \\
$\text{NEREUS}_{\text{best}}$   & 46.85 & 11.15 & 8.26  & 56.02  & 39.84  & 118.58 & 82.35  \\
\bottomrule
\end{tabular}
\end{table}


\begin{figure}[tb]
    \centering
    \includegraphics[width=.9\linewidth]{images/trajectory_grid.pdf}
    \caption{Example predicted trajectories on the FH10 test set. Colour encodes mode probability.}
    \label{fig:example_trajectories}
\end{figure}

% TODO: Insert actual evaluation results table here.
% Suggested structure:
%  - Table 2: ADE / FDE@{1,3,5}min / minADE / minFDE@5min on FH10 test
%  - Table 3: Cross-region results on DMA10
%  - Table 4: Per-ship-type breakdown on FH10
%  - Table 5 / Figure: Collision-risk metrics

\subsection{Module Ablation}

\paragraph{Prior module.}
The prior has the largest and most consistent effect.
Disabling it yields a mean ADE of 59.1\,m,
while any prior variant improves this by roughly 7\,m (density 51.4, path 52.4, cluster 52.6).
The three variants are close to each other, with the density prior slightly ahead,
so the benefit comes mainly from having any prior rather than from the specific formulation.

\paragraph{Social module.}
GAT-based interaction modeling gives the best results (50.4 mean ADE),
while social pooling (56.0) actually performs worse than no social module at all (52.1).
The attention mechanism appears necessary to extract useful interaction information; naive pooling injects noise.

\paragraph{Map module.}
The attention-based map encoder clearly helps (50.4 vs 54.1 without a map),
whereas the CNN encoder is barely distinguishable from no map
(53.6) and shows the highest variance across configurations (std 5.4).
%- its usefulness depends strongly on which other modules are active (e.g. pool+cnn+density is the second-worst configuration overall at 61.4).

\paragraph{Combined.}
The effects compound:
%the best configuration, GAT + attention map + density prior, reaches 46.9 ADE / 35.3 minADE.
%\gaf{das hier mit rein? Notably, off/off/cluster — the cluster prior alone — already ranks fourth overall (48.1 ADE), underscoring the prior module's dominant contribution.}

the best configuration---%
GAT $+$ attention map $+$ density prior---%
reaches $46.9$\,m \ac{ADE} / $35.3$\,m \ac{mADE}.
Notably, the cluster prior alone (no social, no map module) already ranks fourth overall
at $48.1$\,m \ac{ADE}, underscoring the prior module's dominant contribution.

%raus aus platzgründen
% \begin{table}[t]
% \centering
% \caption{Ablation results: effect of each module, reported as mean $\pm$ std over all combinations of the other modules.}
% \label{tab:ablations}
% \begin{tabular}{ll|r|r}
% \toprule
% Module & Variant & ADE & minADE \\
% \midrule
% \multirow{3}{*}{Social} & off  & $52.1 \pm 4.8$ & $40.0 \pm 3.9$ \\
%                         & gat  & $50.4 \pm 3.7$ & $38.5 \pm 3.0$ \\
%                         & pool & $56.0 \pm 3.5$ & $43.0 \pm 2.6$ \\
% \midrule
% \multirow{3}{*}{Map}    & off  & $54.1 \pm 4.6$ & $41.4 \pm 3.5$ \\
%                         & atte & $50.4 \pm 2.7$ & $38.5 \pm 2.4$ \\
%                         & cnn  & $53.6 \pm 5.4$ & $41.3 \pm 4.4$ \\
% \midrule
% \multirow{4}{*}{Prior}  & off     & $59.1 \pm 2.9$ & $44.6 \pm 2.6$ \\
%                         & cluster & $52.6 \pm 5.2$ & $40.4 \pm 4.2$ \\
%                         & density & $51.4 \pm 4.6$ & $39.3 \pm 3.9$ \\
%                         & path    & $52.4 \pm 2.9$ & $40.4 \pm 2.4$ \\
% \bottomrule
% \end{tabular}
% \end{table}

% \begin{figure}[tb]
%     \centering
%     \includegraphics[width=\linewidth]{images/ablation_modules_fde.pdf}
%     \caption{\ac{FDE} versus prediction horizon on the test set, averaged over all submodule
%     combinations for each variant of the social, map, and prior modules.
%     \msp{This figure is in isolation. Would dropping it be problematic?}
%     \gaf{Nah, would be fine.}
%     }
%     \label{fig:fde_horizon}
% \end{figure}


% \begin{figure}
%     \centering
%     \includegraphics[width=0.75\linewidth]{images/shipgroup_modules_k_ade.pdf}
%     \caption{Enter Caption}
%     \label{fig:placeholder}
% \end{figure}
% \placeholder{The progressive addition of modules consistently reduces ADE and
% minFDE@5\,min on the FH10 test set.
% The static module alone (\texttt{base}) outperforms the plain GRU,
% confirming that ship-type information benefits even single-vessel
% prediction.
% Adding the density prior further reduces error, as the model
% leverages historical route distributions.
% The full \texttt{density\_cnn\_gat} variant achieves the lowest
% deterministic ADE, while the map-attention variant \texttt{atte}
% achieves the lowest minADE by effectively querying the relevant scene
% region conditioned on the trajectory prior.}

% \subsection{Cross-Region Generalisation}

% \placeholder{On the DMA10 test set --- unseen geography --- the relative ranking
% of ablation variants is preserved, indicating that each contextual
% module encodes transferable spatial priors rather than region-specific
% memorisation.}

\subsection{Per-Ship-Type Analysis}

%Figure \ref{fig:ade_shipgroups} shows the ADE by Ship-Type and submodule. While there are clear differences in the performance between the ship groups, with cargo ships being much easier to predict and the other category the hardest, there is not strong pattern between Ship-Type and submodule usage.
\Cref{fig:ade_shipgroups} breaks down \ac{ADE} by ship type and submodule.
Performance differs markedly between groups,
\textit{cargo}-vessels are the easiest to predict
and the heterogeneous \textit{other} category the hardest.%
%
Yet no submodule shows a consistent advantage for a particular ship type.
%
% Both observations are consistent with a representation-learning analysis of the same waters~\cite{al2026representation}:
% there,
% \textit{cargo} trajectories form the most spatially regular,
% compact structure (deep-water main-channel transits)
% while
% \textit{other} traffic lacks a coherent behavioural signature,
% and--crucially--learnt behavioural structure is organised by
% \emph{operational context} (channel transit, port manoeuvring, anchorage)
% that cuts across vessel categories rather than by ship type.
% The latter offers a mechanistic explanation for the absence of a ship-type\,$\times$\,module interaction here:
% because the informative axis of variation is operational rather than categorical,
% each contextual module benefits the ship types roughly uniformly instead of selectively.

Both observations match a representation-learning analysis of the same waters \cite{al2026representation},
which found that behavioural structure in these waters is organised by operational context
(channel transit, port manoeuvring, anchorage) rather than ship type,
explaining why each module benefits all ship types roughly uniformly.

\begin{figure}[h]
    \centering
    \includegraphics[width=\linewidth]{images/shipgroup_modules_ade.pdf}
    \caption{\ac{ADE} by submodule and ship type. Boxes aggregate over the other modules.}
    \label{fig:ade_shipgroups}
\end{figure}

% \gaf{n- and m- dashes are used at different locations worng! Check before submittion!}


%            n_samples  n_trajs  mean_speed  sample_share
% sailing      1164063     5400        2.36          31.7
% other        1019677     7280        4.36          27.7
% cargo         870638     3878        4.09          23.7
% passenger     492014     6256        4.25          13.4
% unknown       128242      542         NaN           3.5


% \placeholder{Cargo and passenger vessels, which follow well-defined traffic
% separation lanes, benefit most from the density and cluster prior
% modules.
% Sailing vessels exhibit higher trajectory variance; for this group
% the social GAT module provides the largest improvement, likely because
% sail boats frequently manoeuvre around each other in close proximity.}

% \subsection{Collision Risk Assessment}

% \placeholder{Replacing the context-free \texttt{base} predictions with
% context-aware NEREUS predictions reduces the predicted collision
% ratio across all ship-group strata, suggesting that richer contextual
% information leads to collision-avoiding trajectory hypotheses consistent
% with expected maritime behaviour.}


\paragraph{Limitations.}
While the average forecast is accurate, some situations remain hard.
\Cref{fig:busy_trajectories} shows trajectories from the busiest hour in the dataset:
2023-05-18 was a public holiday that drew a large number of sailing boats onto erratic, wind-driven paths,
on which the constant-heading assumption underlying most modes breaks down.
%
A second, structural limitation concerns spatial generalisation:
because the map and density context are keyed to absolute position,
part of the learned signal is expected to be region-specific rather than transferable.
A systematic cross-region evaluation on geographically distinct waters is therefore left to future work.

\begin{figure}[h]
    \centering
    \includegraphics[width=.9\linewidth]{images/busy_trajectory_grid.pdf}
    \caption{Predicted trajectories from the busiest (public-holiday) day in the dataset,
    dominated by erratically manoeuvring sailing vessels.}
    \label{fig:busy_trajectories}
\end{figure}

% ============================================================
%  7. CONCLUSION
% ============================================================
\section{Conclusion}
\label{sec:conclusion}

We presented NEREUS, a modular, context-aware probabilistic
framework for short-term ship trajectory forecasting.
% Through an ablation study across twelve variants we demonstrated that
%\msp{correct? -> müssen dann noch einen minADE plot rein machen}
Through an ablation across 30~module combinations, we demonstrated that
combining temporal dynamics, static vessel properties, social
interactions modelled by graph attention, and spatial-environmental
map context in a learnable fusion leads to consistent improvements over
single-context baselines.
The three trajectory prior modules--density maps, TrAISformer path
prediction, and GMM cluster maps--provide complementary mechanisms for
incorporating historical route knowledge and improve both deterministic
and stochastic evaluation metrics.

Future work includes extending the graph to include AIS message quality
uncertainty, integrating NEREUS into a real-time collision-avoidance
decision support system, and exploring diffusion-based decoders as an
alternative to the MDN for richer trajectory distributions.

% \begin{credits}
% \subsubsection{\ackname}
% % TODO: add funding acknowledgements.
% \subsubsection{\discintname}
% The authors have no competing interests to declare.
% \end{credits}



% % ============================================================
% %  APPENDIX
% % ============================================================
% \appendix
% %% NOTE: LNCS renders this heading as "A  Input Feature Schema ...". That is the
% %% correct Springer appendix style. If you prefer the word "Appendix" to appear,
% %% uncomment the next line (cleveref-safe; \Cref{app:features} then reads "Appendix A"):
% % \renewcommand{\thesection}{Appendix~\Alph{section}}
% \section{Input Feature Schema and Preprocessing}
% \label{app:features}

% Raw \ac{AIS} from multiple sources is decoded, outlier-filtered, segmented at
% transmission gaps and stationary episodes, Kalman-smoothed, and interpolated onto
% a fixed 10\,s grid in UTM coordinates (EPSG:25832); the full pipeline and its
% validation are reported in~\cite{al2026representation}.
% \Cref{tab:features} lists the features the model consumes, in code order.
% Angular quantities (course, turn rate, course difference, relative bearing) are
% sine/cosine encoded; edge distances (distance, \ac{TCPA}, \ac{DCPA}) are
% $\log(1{+}\cdot)$-scaled; the collision-risk score is passed through unnormalised.

% \begin{table}[htb]
% \caption{NEREUS input features by group, in code order.
% Node dim 8; static dim 8; edge dim 25; four map channels; single-channel prior.}
% \label{tab:features}
% \centering\small
% \setlength{\tabcolsep}{6pt}
% \begin{tabular}{p{2.5cm}p{8.5cm}}
% \toprule
% Group & Features (per time step unless static), in order \\
% \midrule
% Trajectory (8) & $\Delta x$; $\Delta y$; SOG; acceleration;
% $\sin$/$\cos$ course; $\sin$/$\cos$ turn rate \\
% Static (8) & hull offsets \{to\_bow, to\_stern, to\_port, to\_starboard\};
% ship-type one-hot \{sailing, cargo, passenger, other\} \\
% Interaction (25) & distance; relative SOG; \acs{TCPA}; \acs{DCPA};
% collision-risk; $\sin$/$\cos$ course-difference; $\sin$/$\cos$ relative bearing;
% ego static (8); neighbour static (8) \\
% Map (4) & distance to coastline; distance to restricted area; distance to ferry
% route; water depth ($500\,$m crop, 50\,m grid) \\
% Prior (1) & ship-type density map \emph{or} TrAISformer-path heat-map \emph{or}
% \acs{GMM} cluster-density mixture \\
% \bottomrule
% \end{tabular}
% \end{table}

% ============================================================
%  BIBLIOGRAPHY
% ============================================================
\bibliographystyle{preamp/splncs04}
\bibliography{bibliography}

\end{document}